\documentclass[../../../include/open-logic-section]{subfiles}

\begin{document}

\olfileid{sth}{ordinals}{ordtype}
\olsection{క్రమరకాలుగా క్రమసంఖ్యలు}

ప్రతిస్థాపనతో, అంటే
ఇప్పుడు $\ZFminus$లో
పని చేస్తూ, మనం
లక్ష్యంగా పెట్టుకున్న
ఫలితాన్ని చివరకు
నిరూపించవచ్చు:

\begin{thm}\ollabel{thmOrdinalRepresentation}
ప్రతి సుక్రమమూ ఒక
ఏకైక క్రమసంఖ్యతో
సమరూపమవుతుంది.
\end{thm}

\begin{proof}
$\tuple{A, <}$ ఒక
సుక్రమం అనుకుందాం.
\olref[basic]{ordisoidentity}
ప్రకారం, అది గరిష్ఠంగా
ఒక క్రమసంఖ్యతో
సమరూపమవుతుంది.
కాబట్టి వైరుధ్యాన్ని
చూపడానికి, $\tuple{A, <}$
\emph{ఏ} క్రమసంఖ్యతోనూ
సమరూపం కాదనుకుందాం.
మొదట $\tuple{A, <}$ను
``సాధ్యమైనంత చిన్నదిగా''
చేస్తాం. వివరంగా:
దాని తనకన్నా చిన్న
ఆరంభ ఖండం
$\tuple{A_a,
<_a}$ ఏ క్రమసంఖ్యతోనూ
సమరూపం కాకపోతే,
ఆ ధర్మం గల అత్యల్ప
$a \in A$ ఉంటుంది;
అప్పుడు $B = A_a$,
$\mathord{\lessdot} =
\mathord{<_a}$ అనుకుందాం.
లేకపోతే $B = A$,
$\mathord{\lessdot} =
\mathord{<}$ అనుకుందాం.

నిర్వచనం వల్ల $B$
యొక్క ప్రతి తనకన్నా
చిన్న ఆరంభ ఖండం
ఏదో ఒక క్రమసంఖ్యతో
సమరూపం; అది పైన
చెప్పినట్టు ఏకైకం.
కాబట్టి ప్రతిస్థాపన
ప్రకారం ఈ కింది సమితి
ఉంటుంది; అది ఒక
ప్రమేయం కూడా:
\[
	f = \Setabs{\tuple{\beta, b}}{b \in B\text{ మరియు }\ordeq{\beta}{\tuple{B_b, \lessdot_b}}}
\]
వైరుధ్యాన్ని పూర్తి
చేయడానికి, ఏదో
క్రమసంఖ్య $\alpha$కు
$f$ ఒక $\alpha \to B$
సమరూపత అని చూపుతాం.

$\ran{f}
= B$ అనేది స్పష్టం.
\olref[iso]{lemordsegments}
ప్రకారం $f$ క్రమాన్ని
సంరక్షిస్తుంది; అంటే,
$\gamma \in \beta$
అప్పుడే మరియు అప్పుడే
$f(\gamma) \lessdot f(\beta)$.
$\dom{f}$ క్రమసంఖ్య
అని చూపడానికి,
\olref[basic]{corordtransitiveord}
ప్రకారం $\dom{f}$
సంక్రామకమని చూపితే
సరిపోతుంది. అందుకోసం
$\beta \in \dom{f}$ను
స్థిరపరచండి; అంటే,
ఏదో $b$కు
$\ordeq{\beta}{\tuple{B_b, \lessdot_b}}$.
$\gamma \in
\beta$ అయితే,
\olref[iso]{wellordinitialsegment}
ప్రకారం
$\gamma \in \dom{f}$;
సాధారణీకరిస్తే,
$\beta \subseteq
\dom{f}$.
\end{proof}

ఈ ఫలితంతో,
\olref[vn]{sec} నుంచి
ఇవ్వాలనుకున్న ఈ
నిర్వచనాన్ని ఇప్పుడు
ఇవ్వవచ్చు:

\begin{defn}
$\tuple{A, <}$ ఒక
సుక్రమం అయితే,
దాని క్రమరకం
$\ordtype{A, <}$ అంటే
$\ordeq{\tuple{A, < }}{\alpha}$
అయ్యే ఏకైక
క్రమసంఖ్య $\alpha$.
\end{defn}

ఇంకా ఈ నిర్వచనం
రెండు మంచి సూత్రాలను
ఇస్తుంది:

\begin{cor}\ollabel{ordtypesworklikeyouwant}
$\tuple{A, <}$,
$\tuple{B, \lessdot}$
సుక్రమాలైతే:
\begin{align*}
	\ordtype{A, <} = \ordtype{B, \lessdot}&\text{ అప్పుడే మరియు అప్పుడే }\ordeq{\tuple{A, <}}{\tuple{B, \lessdot}}\\
	\ordtype{A, <} \in \ordtype{B, \lessdot}&\text{ అప్పుడే మరియు అప్పుడే }\ordeq{\tuple{A, <}}{\tuple{B_b, \lessdot_b}}\text{ ఏదో ఒక }b \in B
\end{align*}
\end{cor}

\begin{proof}
సమానత్వం
\olref[basic]{ordisoidentity}
ప్రకారం వస్తుంది.
రెండవ వాదన కోసం
$\ordtype{A, <} = \alpha$,
$\ordtype{B,
\lessdot} = \beta$
అనుకుందాం; మన సమరూపత
$f \colon \beta \to B$
అనుకుందాం.
\sourcecorrection{OLTESTORDTYPE-001}{మూలం సమరూపత ప్రమేయానికి క్రమయుగ్మాన్నే లక్ష్య సమితిగా రాసింది; ఇక్కడ దాని సమితి లక్ష్యం అని సరిచేశాం.}
అప్పుడు:
\begin{align*}
	\alpha \in \beta&\text{ అప్పుడే మరియు అప్పుడే }\ordeq{\alpha}{\tuple{B_b, \lessdot_b}}\text{ ఏదో ఒక }b \in B\\
	&\text{ అప్పుడే మరియు అప్పుడే }\ordeq{\tuple{A, <}}{\tuple{B_b, \lessdot_b}}\text{ ఏదో ఒక $b \in B$ కోసం}
\end{align*}
\sourcecorrection{OLTESTORDTYPE-002}{మూల సమానార్థక శ్రేణిలో సభ్యత్వం అసత్యమైతే నిర్వచితం కాని ప్రమేయ విలువను వాడారు; ఇక్కడ ఏదో ఒక ఆరంభ ఖండంతో సమరూపత రూపంలో రెండు దిశలనూ ప్రకటించాం.}
ఈ సమానార్థకాలు
\olref[iso]{ordisounique},
\olref[iso]{wellordinitialsegment},
\olref[basic]{ordissetofsmallerord}
ద్వారా వస్తాయి.
\end{proof}

\end{document}
